% Bölüm 17 — taslak plan. Tam metin yazıldığında bu dosyanın yerini alır.
\chapter{Kip Mantığı ve Kripke Anlambilimi}\label{bol:kip}

\begin{bolumhedefleri}
\begin{itemize}
  \item Zorunluluk ({\sembolfont □}) ve olanak ({\sembolfont ◇}) işlemlerini
  \item Kripke çerçevelerini ve modellerini
  \item K, T, S4, S5 sistemlerini ve çerçeve koşullarını
  \item Kip mantığının Lean'de formalizasyonunu
\end{itemize}
\end{bolumhedefleri}

\begin{taslakplan}
\textbf{Planlanan kesimler}
\begin{enumerate}[label=\thechapter.\arabic*, leftmargin=2.4em]
  \item Kip dili: {\sembolfont □} ve {\sembolfont ◇}; okunuşlar (zorunlu, bilinen, kanıtlanabilir)
  \item Kripke çerçeveleri
  \item Karşılık kuramı: T $↔$ yansımalı, 4 $↔$ geçişli
  \item Kanonik modeller ve tamlık
  \item Foundation projesinde kip mantığı
\end{enumerate}

\textbf{Planlanan Lean örnekleri}
\begin{itemize}
  \item \lc{Kripke modeli tanımı}
  \item \lc{□(p → q) → □p → □q (K aksiyomu) geçerliliği}
  \item \lc{Yansımalı çerçevede □p → p}
\end{itemize}

\textbf{Temel kaynaklar:} \citet{kripke1963,blackburn2001,foundation}.
\end{taslakplan}

